% Блок-схема функции parser::tokenizer::tokenize() — лексический анализ
% ГОСТ 19.701-90, src/parser/tokenizer.cpp

\documentclass[a4paper]{article}
\usepackage[T2A]{fontenc}
\usepackage[utf8]{inputenc}
\usepackage[russian]{babel}
\usepackage[a4paper,margin=1.5cm]{geometry}
\usepackage{tikz}
\usetikzlibrary{shapes.geometric, arrows.meta, positioning, calc}

\tikzset{
  terminal/.style={
    rectangle, rounded corners=8pt,
    draw=black, line width=1pt,
    minimum width=3.8cm, minimum height=0.9cm,
    text centered, font=\small
  },
  process/.style={
    rectangle,
    draw=black, line width=1pt,
    minimum width=3.8cm, minimum height=0.9cm,
    text centered, font=\small
  },
  decision/.style={
    diamond, aspect=2.6,
    draw=black, line width=1pt,
    inner sep=1pt,
    minimum width=3cm, minimum height=0.9cm,
    text centered, font=\footnotesize
  },
  predefined/.style={
    rectangle, double, double distance=2pt,
    draw=black, line width=1pt,
    minimum width=3.8cm, minimum height=0.9cm,
    text centered, font=\small
  },
  arrow/.style={->, >=Stealth, line width=1pt}
}

\begin{document}
\begin{center}
\textbf{Блок-схема функции \texttt{tokenize(expression)}}
\end{center}

\begin{center}
\begin{tikzpicture}[node distance=0.6cm]

\node (start)  [terminal]                        {Начало tokenize};
\node (init)   [process,   below=of start]       {tokens = []; \quad $i = 0$; \quad $n = |\text{expr}|$};

\node (cond)   [decision,  below=of init]        {$i < n$?};
\node (post)   [decision,  right=3cm of cond]    {tokens пуст?};
\node (errE)   [terminal,  below=of post]        {Ошибка: пустое выражение};
\node (ret)    [terminal,  right=2cm of post]    {Возврат tokens};

\node (c)      [process,   below=of cond]        {$c = $ expr[$i$]};
\node (sp)     [decision,  below=of c]           {$c$ — пробел?};
\node (incSp)  [process,   right=3cm of sp]      {$i = i + 1$};

\node (dig)    [decision,  below=of sp]          {$c$ — цифра или '.'?};
\node (rdNum)  [predefined,right=3cm of dig]     {read\_number $\to$ tokens};

\node (al)     [decision,  below=of dig]         {$c$ — буква?};
\node (rdFun)  [predefined,right=3cm of al]      {read\_function $\to$ tokens};

\node (lp)     [decision,  below=of al]          {$c$ = '(' или ')'?};
\node (addP)   [process,   right=3cm of lp]      {push скобку; \quad $i = i + 1$};

\node (op)     [decision,  below=of lp]          {$c \in \{+, -, *, /\}$?};
\node (rdOp)   [predefined,right=3cm of op]      {read\_operator $\to$ tokens};

\node (errU)   [terminal,  below=of op]          {Ошибка: неизвестный символ};

\draw [arrow] (start) -- (init);
\draw [arrow] (init)  -- (cond);

\draw [arrow] (cond)  -- node[above] {Нет} (post);
\draw [arrow] (post)  -- node[left]  {Да}  (errE);
\draw [arrow] (post)  -- node[above] {Нет} (ret);

\draw [arrow] (cond)  -- node[left]  {Да}  (c);
\draw [arrow] (c)     -- (sp);

\draw [arrow] (sp)    -- node[above] {Да} (incSp);
\draw [arrow] (incSp) -- ++(3.2,0) |- (cond.east);

\draw [arrow] (sp)    -- node[left]  {Нет} (dig);
\draw [arrow] (dig)   -- node[above] {Да}  (rdNum);
\draw [arrow] (rdNum) -- ++(3.2,0) |- (cond.east);

\draw [arrow] (dig)   -- node[left]  {Нет} (al);
\draw [arrow] (al)    -- node[above] {Да}  (rdFun);
\draw [arrow] (rdFun) -- ++(3.2,0) |- (cond.east);

\draw [arrow] (al)    -- node[left]  {Нет} (lp);
\draw [arrow] (lp)    -- node[above] {Да}  (addP);
\draw [arrow] (addP)  -- ++(3.2,0) |- (cond.east);

\draw [arrow] (lp)    -- node[left]  {Нет} (op);
\draw [arrow] (op)    -- node[above] {Да}  (rdOp);
\draw [arrow] (rdOp)  -- ++(3.2,0) |- (cond.east);

\draw [arrow] (op)    -- node[left]  {Нет} (errU);

\end{tikzpicture}
\end{center}

\end{document}
